\section{Drift, heading and the compass}
\label{sec:corr-heading}

Current and wind carry the boat over the ground. Heading enters when world
velocity is expressed in the body frame, so heading errors affect both apparent
cross response and the drift used by the controller.

\subsection{World drift and body-frame response}

The conditions estimator pools neutral-window world velocity to measure drift
$\dw$. The controller rotates this drift into the body frame as
$R(\psi)^\top\dw$ (Section~\ref{sec:controller}). The offline characterization
response instead subtracts the preceding neutral's mean body velocity from the
hold's mean body velocity; it uses the recorded odometry components directly,
without antenna lever-arm correction or reconstructed hold heading.

\subsection{What a heading error does}

Let the heading used be $\psi + \delta$ instead of $\psi$. Two things go wrong
(Figure~\ref{fig:corr-heading}).

\paragraph{The drift points the wrong way.} The body-frame drift becomes
$R(\psi+\delta)^\top\dw = R(\psi)^\top R(\delta)^\top\dw$: the right vector turned
by $-\delta$. The error is $R(\psi)^\top(R(\delta)^\top - I)\dw$, and since a
rotation keeps lengths, its size is $|(R(\delta) - I)\dw|$. The columns of
$R(\delta) - I$ have length $\sqrt{(1-\cos\delta)^2 + \sin^2\delta} =
\sqrt{2 - 2\cos\delta} = 2|\sin(\delta/2)|$, so
\[
  |\text{drift error}| = 2\,|\dw|\,|\sin(\delta/2)| \approx |\dw|\,|\delta| .
\]
A $10^\circ$ error misplaces $17\,\%$ of the drift. A constant $\delta$ turns the
drift arrow; a heading-dependent error (compass deviation) makes it wobble as the
boat turns.

\paragraph{Straight motion looks like crab.} A boat moving straight ahead at $V$
has world velocity $R(\psi)(V, 0)$; turned back with the wrong heading,
\[
  R(\psi+\delta)^\top R(\psi)\begin{bmatrix}V\\0\end{bmatrix}
  = R(\delta)^\top\begin{bmatrix}V\\0\end{bmatrix}
  = \begin{bmatrix}V\cos\delta\\ -V\sin\delta\end{bmatrix},
\]
using $R(a)^\top R(b) = R(b - a)$. The apparent crab is
\[
  \frac{\text{cross}}{\text{main}} = -\tan\delta ,
\]
approximately the negative heading error, in radians. Against the crab tolerance of $0.05$ that is
large: $1.5^\circ$ already uses half of it, and a compass is good to a few degrees
at best. This rotation can therefore dominate a small measured cross response.

\begin{figure}[htbp]
\centering
\begin{tikzpicture}[scale=2.4, >=Stealth]
  \begin{scope}
    % d at 30 degrees, turned by -20 degrees
    \draw[->, thick, sC] (0,0) -- ({1.5*cos(30)},{1.5*sin(30)}) node[above, font=\small] {$\dw$};
    \draw[->, thick, sB] (0,0) -- ({1.5*cos(10)},{1.5*sin(10)}) node[below right, font=\small] {$R(\delta)^\top\dw$};
    \draw[sD] ({0.55*cos(30)},{0.55*sin(30)}) arc (30:10:0.55);
    \node[font=\scriptsize] at ({0.68*cos(20)},{0.68*sin(20)}) {$\delta$};
    \draw[sA, very thick] ({1.5*cos(30)},{1.5*sin(30)}) -- ({1.5*cos(10)},{1.5*sin(10)});
    \node[sA, font=\scriptsize, anchor=west] at ({1.5*cos(20)+0.08},{1.5*sin(20)+0.05}) {$2|\dw|\sin\frac{\delta}{2}$};
    \node[font=\small] at (0.75,-0.45) {drift turned by $\delta$};
  \end{scope}
  \begin{scope}[xshift=2.7cm]
    \draw[->, thick, sC] (0,0) -- (1.5,0) node[above right, font=\small] {true $(V,0)$};
    \draw[->, thick, sB] (0,0) -- ({1.5*cos(-20)},{1.5*sin(-20)}) node[right, font=\small] {seen};
    \draw[sA, very thick, dashed] ({1.5*cos(-20)},0) -- ({1.5*cos(-20)},{1.5*sin(-20)});
    \node[sA, font=\scriptsize, anchor=east] at ({1.5*cos(-20)-0.03},{0.75*sin(-20)}) {$V\sin\delta$};
    \node[font=\small] at (0.75,-0.75) {straight motion seen as crab};
  \end{scope}
\end{tikzpicture}
\caption{A heading error $\delta$ (here $20^\circ$) turns the body-frame drift
(left) and turns straight motion into an apparent sideways velocity (right). The
drift error does not grow with speed; the fake crab does.}
\label{fig:corr-heading}
\end{figure}

\paragraph{When heading is unreliable.} The diagnostic flags compass--gyro
change disagreement above $5^\circ$ over rolling 10\,s windows, magnetometer
calibration residual above $10^\circ$, yaw standard deviation at least $0.5$\,rad,
or field magnitude outside $0.7$--$1.3$ times the expected field.

\subsection{The thrusters' magnetic field}
\label{sec:corr-mag}

\paragraph{The sequence.} With the props off or the boat tied up, pulse each thruster
alone to $\pm25$, $\pm50$ and $\pm75\,\%$ for 2\,s, with 2\,s of neutral between
pulses and neutral before and after. Record the magnetometer's field in the body
frame (back in mG from the FCU's raw readings).

\paragraph{The model.} The disturbance $\Delta\bm B = (\Delta B_x, \Delta B_y)$ is
the field during a pulse minus the neutral baseline (the mean of the neutral before
and after, which removes slow drift linearly). It is modelled per field axis as
\[
  \Delta B_b = \sum_{i}\Big(\mu^+_{i,b}\,\max(t_i, 0) + \mu^-_{i,b}\,\max(-t_i, 0)\Big),
  \qquad b\in\{x, y\},
\]
eight coefficients per axis: each thruster, each direction, because the current path
and the motor's own field need not be symmetric. It is linear in the $\mu$'s, so it
is least squares once more (Section~\ref{sec:char}), and because each pulse drives one
thruster in one direction, each column of the design matrix is nonzero only during
its own pulses: the columns are orthogonal when both signs of every thruster are sampled. The
coefficients are reported in mG per full effort.

\paragraph{From field to heading error.} The diagnostic uses the horizontal
reference field $B_h=235$\,mG and reports
\[
  \epsilon=\operatorname{atan2}(D,B_h),\qquad
  D=\begin{cases}|\Delta B_\perp|,&\text{with a measured field direction},\\
                |\Delta\bm B|,&\text{without one}.\end{cases}
\]
The latter is the reported heading-independent estimate, not an exact maximum
of the vector-sum angle. It gives $2.4^\circ$ for $10$\,mG and $9.7^\circ$ for
$40$\,mG. Under $3^\circ$ the fit is accepted; $3$--$10^\circ$ requests rerouting
the leads; above $10^\circ$ requests moving the compass.

\begin{figure}[htbp]
\centering
\begin{tikzpicture}
\begin{axis}[corr axis, width=0.62\linewidth, height=5cm, xlabel={thruster effort $t_i$},
  ylabel={$\Delta B_\perp$ (mG)}, xmin=-1, xmax=1, ymin=-30, ymax=40,
  legend pos=north west, legend style={font=\scriptsize}, legend cell align=left]
  \addplot[only marks, mark=*, mark size=1.3pt, sD] table[corr, x index=0, y index=1] {\corrdata{mag}};
  \addlegendentry{pulses}
  \addplot[sA, thick, domain=0:1] {35.3*x};
  \addlegendentry{$\mu^+ = 35$\,mG ($8.5^\circ$ at full)}
  \addplot[sB, thick, domain=-1:0] {22.5*x};
  \addlegendentry{$\mu^- = -22$\,mG ($-5.5^\circ$ at full reverse)}
\end{axis}
\end{tikzpicture}
\caption{The interference fit for one thruster on synthetic pulses: the field
change is fitted separately for each direction, and converted to a heading error
with $\arctan(\Delta B_\perp/235\,\text{mG})$. Data:
\texttt{figures/correction/correction\_tracking.py}.}
\label{fig:corr-mag}
\end{figure}

\file{mhseals\_testing/procedures/characterize.py (response); analysis/conditions.py; analysis/thrust.py (body\_drift); procedures/interference.py; analysis/interference.py (fit, heading\_error); analysis/magnetic.py (compass\_gyro, yaw\_quality)}
